%% LyX 2.3.6 created this file.  For more info, see http://www.lyx.org/.
%% Do not edit unless you really know what you are doing.
\documentclass[a4paper,english]{scrartcl}
\usepackage[T1]{fontenc}
\usepackage[latin9]{inputenc}
\usepackage{amsbsy}

\makeatletter

%%%%%%%%%%%%%%%%%%%%%%%%%%%%%% LyX specific LaTeX commands.
\pdfpageheight\paperheight
\pdfpagewidth\paperwidth


\makeatother

\usepackage{babel}
\usepackage{listings}
\renewcommand{\lstlistingname}{Listing}

\begin{document}
\title{SEA3004F M3P2: Masking, divergence, and curl }
\date{Ocean and Atmosphere Dynamics}
\maketitle

\section{Tutorial 1: array masking [10]}

Download all files and put them in a folder on your computer. Open
Matlab and navigate to that folder.

Fields in earth sciences may not have data at all grid points. Missing
values can represent land points (a mask), lack of measurements or
information and non-representable numbers coming from numerical errors.
The latter case was highlighted in the previous practical, in which
the computation of the numerical gradient with values close to 0 led
to indefinite numbers. It is often sufficient to change the ranges
and/or color scales of the visualization, but a more effective solution
is to mask the data.

Masked data are usually identified with a special number, which in
most computing software is called Not-A-Number (\texttt{NaN}).
In this tutorial, we will learn how to replace certain values of the
field with the missing value NaN (a method called ``logical indexing'').

File \texttt{NCEP\_wind\_matlab.mat} contains the global climatological
wind field (u and v components) from the NCEP atmospheric reanalyses
(1981-2010). This dataset contains two types of coordinates (geographical
and geodetic), as well as an additional array that contains the land
mask. 

We will first compute the wind speed and then mask the land points
using the provided array \texttt{land}. We will then plot the same
wind speed field using first latitude/longitude and then distances
on the sphere (using the same arrays analysed in the previous tutorial).

\begin{lstlisting}[language=Matlab,basicstyle={\small},showstringspaces=false,frame=lines]
load NCEP_wind_matlab.mat
% first visualize the land array using pcolor

<your command here>

% compute the wind speed
wndsp=sqrt(uwnd.^2+vwnd.^2);
% create a mask containing land data
wndsp(land==1)=NaN; % What type of array is this one?

% Plot wind speed using 2 different coordinate systems 
figure
subplot(2,1,1)
pcolor(lon,lat,wndsp)
axis tight
colorbar
shading flat
grid on
title('NCEP Wind speed annual mean (m/s)')
xlabel('Longitude')
ylabel('Latitude')
subplot(2,1,2)
pcolor(Xkm,Ykm,wndsp)
axis tight
colorbar
title('NCEP Wind speed annual mean (m/s)')
xlabel('Distance (km)')
ylabel('Distance (km)')
\end{lstlisting}


\section{Tutorial 2: curl of a vector field [10]}

In this tutorial we will plot an analytical vector field, which represents
some circulation patterns in geophysical fluid dynamics, and we will
compute the vertical component of the curl (sometimes indicated as
$\mathsf{curl}_{z}$)
\[
\boldsymbol{\hat{k}}\cdot\boldsymbol{\nabla\times}\mathbf{F}=\frac{\partial F_{y}}{\partial x}-\frac{\partial F_{x}}{\partial y}
\]

The vector field is 
\[
\mathbf{F}\left(x,y\right)=\sin\left(x+2y\right)\hat{\mathbf{i}}+\cos\left(x-2y\right)\hat{\mathbf{j}}
\]
where $x\in\left[-2,2\right]$ and $y\in\left[-2,2\right]$ in discrete
steps of 0.25. Notice the use of function \texttt{colorposneg} to
produce a positive/negative colorbar.

\begin{lstlisting}[language=Matlab,basicstyle={\small},showstringspaces=false,frame=lines]
% First create the axes
x=-2:.25:2;
y=-2:.25:2;
[X,Y]=meshgrid(x,y);
% Define the components of the vector field
Fx=sin(X+2*Y);
Fy=cos(X-2*Y);

% Graph the vector field using the quiver function
figure
quiver(X,Y,Fx,Fy,'k')
title('Velocity field')
axis equal
axis tight

% Compute, graph the curl and overlay the vectors
figure
C = curl(X,Y,Fx,Fy);
pcolor(X,Y,C)
axis equal
axis tight
hold
quiver(X,Y,Fx,Fy,'k')
title('Velocity field and curl')
colorposneg
colorbar

% optional for printing the figure
print -dpng vector_field_curl.png
\end{lstlisting}


\section{Exercise 1: divergence of a vector field {[}20{]}}

Given the vector field $\mathbf{F}\left(x,y\right)=-2x\hat{\mathbf{i}}-2y\hat{\mathbf{j}}$
, where $x\in\left[-10,10\right]$ and $y\in\left[-10,10\right]$.
\begin{enumerate}
\item {[}10{]} Compute the divergence $\boldsymbol{\nabla}\cdot\mathbf{F}$
analytically (you can use the equation editor of the word processor
or include a picture of your calculations on paper) .
\item {[}10{]} Compare the result above with the one obtained using the function \texttt{divergence}. Carefully read the help on how
to use it. Plot it graphically and overlay the vector field; save
the figure.
\end{enumerate}

\section{Exercise 2: wind stress curl {[}40{]}}

Use the wind data in file \texttt{NCEP\_wind\_matlab.mat} to calculate
and plot the climatological annual wind stress curl. You will have
to create the following two separate figures
\begin{enumerate}
\item A \textbf{masked} latitude/longitude map of wind stress, first computing the wind stress components (use a value of the drag coefficient $C_{D}=1.5\:10^{-3}$) and then the intensity as done in Tutorial 1 above. Add all the labels and the colorbar. Indicate the proper units in the title. [15]
\item A \textbf{masked} latitude/longitude map of the wind stress curl, using the wind stress components calculated above and Tutorial 2 for guidelines on how to compute the curl. Remember that distances must be in m. Add all the labels and the colorbar. Indicate the proper
units in the title. In the caption to this figure, briefly illustrate
which oceanographic regions show the largest/lowest values of the
curl. [25]
\end{enumerate}

\end{document}
